% questions.tex —— 2014 年·山东卷·理科 · 题目内容片段（无答案）
% 由 tools/generate_projects.py 依来源总账自 2014/2014_数学真题_山东理_试题.md 生成，勿手改（改源 MD 后重跑）
\begin{gknotice}
\textbf{来源说明}：本题面依据固定源 PDF 的两页 A3 题面逐题转录，原卷印刷页码为 907—908。第 907 面为选择题 1—10、填空题 11—15 与解答题 16—17；第 908 面为解答题 18—21。全卷三处节头依次为「一、选择题」「二、填空题」「三、解答题」。

\textbf{解析归属}：固定源不附答案或解析；配套答案与解析均由本库按本题面独立复算。

\textbf{图形说明}：第 7 题的频率分布直方图、第 11 题的程序框图、第 17 题的四棱柱、第 18 题的乒乓球台面分区图由项目内 TikZ 矢量源重绘；其余各题原页无图。
\end{gknotice}
\gkpart{一、选择题}
\begin{question}
\gkti{1}已知 $a,b\in\mathbf R$，$\mathrm i$ 是虚数单位，若 $a-\mathrm i$ 与 $2+b\mathrm i$ 互为共轭复数，则 $(a+b\mathrm i)^2=$\mbox{（　　）}
（A）$5-4\mathrm i$；（B）$5+4\mathrm i$；（C）$3-4\mathrm i$；（D）$3+4\mathrm i$。
\end{question}
\begin{question}
\gkti{2}设集合 $A=\{x\mid |x-1|<2\}$，$B=\{y\mid y=2^x,\ x\in[0,2]\}$，则 $A\cap B=$\mbox{（　　）}
（A）$[0,2]$；（B）$(1,3)$；（C）$[1,3)$；（D）$(1,4)$。
\end{question}
\begin{question}
\gkti{3}函数 $f(x)=\dfrac1{\sqrt{(\log_2x)^2-1}}$ 的定义域为\mbox{（　　）}
（A）$\left(0,\dfrac12\right)$；（B）$(2,+\infty)$；（C）$\left(0,\dfrac12\right)\cup(2,+\infty)$；（D）$\left(0,\dfrac12\right]\cup[2,+\infty)$。
\end{question}
\begin{question}
\gkti{4}用反证法证明命题“设 $a,b$ 为实数，则方程 $x^3+ax+b=0$ 至少有一个实根”时，要做的假设是\mbox{（　　）}
（A）方程 $x^3+ax+b=0$ 没有实根；（B）方程 $x^3+ax+b=0$ 至多有一个实根；（C）方程 $x^3+ax+b=0$ 至多有两个实根；（D）方程 $x^3+ax+b=0$ 恰好有两个实根。
\end{question}
\begin{question}
\gkti{5}已知实数 $x,y$ 满足 $a^x<a^y$（$0<a<1$），则下列关系式恒成立的是\mbox{（　　）}
（A）$\dfrac1{x^2+1}>\dfrac1{y^2+1}$；（B）$\ln(x^2+1)>\ln(y^2+1)$；（C）$\sin x>\sin y$；（D）$x^3>y^3$。
\end{question}
\begin{question}
\gkti{6}直线 $y=4x$ 与曲线 $y=x^3$ 在第一象限内围成的封闭图形的面积为\mbox{（　　）}
（A）$2\sqrt2$；（B）$4\sqrt2$；（C）$2$；（D）$4$。
\Needspace{0.4\textheight}
\end{question}
\begin{question}
\gkti{7}为了研究某种药品的疗效，选取若干名志愿者进行临床试验。所有志愿者的舒张压数据（单位：kPa）的分组区间为 $[12,13)$，$[13,14)$，$[14,15)$，$[15,16)$，$[16,17]$，将其按从左到右的顺序分别编号为第一组，第二组，…，第五组。如图所示是根据试验数据制成的频率分布直方图。已知第一组与第二组共有 $20$ 人，第三组中没有疗效的有 $6$ 人，则第三组中有疗效的人数为\mbox{（　　）}
（A）$6$；（B）$8$；（C）$12$；（D）$18$。
\nopagebreak[4]
\begin{center}% ============================================================
% f2014_shandong_science_q07_histogram.tex
% 2014 山东卷（理科）第 7 题：舒张压数据的频率分布直方图
% 来源 PDF：普通高考/2014/2014山东理.pdf
% 原页：第 1 页（A3 横排）左栏，印面 907，第 7 题题干下方
% 各矩形高（频率/组距）：0.24、0.16、0.36、0.16、0.08，和为 1。
% ============================================================
\begin{tikzpicture}[>=Stealth, line width=0.7pt, font=\small,
  x=1cm, y=1cm]
  % 坐标轴
  \draw[->] (0.60,0) -- (7.35,0) node[below right] {舒张压/kPa};
  \draw[->] (0.60,-0.05) -- (0.60,4.05);
  \node[below left] at (0.60,0) {$0$};
  % 纵轴标题：频率 / 组距（上下两行，原图无分数线）
  \node at (0.62,4.62) {频率};
  \node at (0.62,4.16) {组距};
  % x 轴断开符号
  \draw (0.86,0) -- (0.98,0.10) -- (1.10,-0.10) -- (1.22,0.10) -- (1.34,0);
  % 直方图：组距 1，映射 12→1.90，13→2.90，…，17→6.90
  \draw (1.90,0) rectangle (2.90,2.40);
  \draw (2.90,0) rectangle (3.90,1.60);
  \draw (3.90,0) rectangle (4.90,3.60);
  \draw (4.90,0) rectangle (5.90,1.60);
  \draw (5.90,0) rectangle (6.90,0.80);
  % 纵轴刻度与虚线
  \foreach \h/\y in {0.08/0.80, 0.16/1.60, 0.24/2.40, 0.36/3.60}
  {
    \draw (0.54,\y) -- (0.66,\y);
    \node[left] at (0.52,\y) {$\h$};
    \draw[dashed] (0.66,\y) -- (6.90,\y);
  }
  % 横轴刻度
  \foreach \x/\l in {1.90/12, 2.90/13, 3.90/14, 4.90/15, 5.90/16, 6.90/17}
  {
    \draw (\x,0.05) -- (\x,-0.05);
    \node[below] at (\x,-0.02) {$\l$};
  }
\end{tikzpicture}
\end{center}
\nopagebreak[4]
\end{question}
\begin{question}
\gkti{8}已知函数 $f(x)=|x-2|+1$，$g(x)=kx$，若方程 $f(x)=g(x)$ 有两个不相等的实根，则实数 $k$ 的取值范围是\mbox{（　　）}
（A）$\left(0,\dfrac12\right)$；（B）$\left(\dfrac12,1\right)$；（C）$(1,2)$；（D）$(2,+\infty)$。
\end{question}
\begin{question}
\gkti{9}已知 $x,y$ 满足约束条件 $\begin{cases}x-y-1\leqslant0,\\2x-y-3\geqslant0,\end{cases}$ 当目标函数 $z=ax+by$（$a>0$，$b>0$）在该约束条件下取到最小值 $2\sqrt5$ 时，$a^2+b^2$ 的最小值为\mbox{（　　）}
（A）$5$；（B）$4$；（C）$\sqrt5$；（D）$2$。
\end{question}
\begin{question}
\gkti{10}已知 $a>b>0$，椭圆 $C_1$ 的方程为 $\dfrac{x^2}{a^2}+\dfrac{y^2}{b^2}=1$，双曲线 $C_2$ 的方程为 $\dfrac{x^2}{a^2}-\dfrac{y^2}{b^2}=1$，$C_1$ 与 $C_2$ 的离心率之积为 $\dfrac{\sqrt3}2$，则 $C_2$ 的渐近线方程为\mbox{（　　）}
（A）$x\pm\sqrt2y=0$；（B）$\sqrt2x\pm y=0$；（C）$x\pm2y=0$；（D）$2x\pm y=0$。
\end{question}
\gkpart{二、填空题}
\Needspace{0.4\textheight}
\begin{question}
\gkti{11}执行如图所示的程序框图，若输入的 $x$ 的值为 $1$，则输出的 $n$ 的值为\gkblank{}。
\nopagebreak[4]
\begin{center}% ============================================================
% f2014_shandong_science_q11_flowchart.tex
% 2014 山东卷（理科）第 11 题：程序框图
% 来源 PDF：普通高考/2014/2014山东理.pdf
% 原页：第 1 页（A3 横排）中栏，印面 907，第 11 题题干下方
% 结构：开始 → 输入 x → n=0 → 判断 x^2-4x+3≤0；
%       是：x=x+1 → n=n+1 → 返回判断；否：输出 n → 结束。
% ============================================================
\begin{tikzpicture}[>=Stealth, line width=0.7pt, font=\small,
  x=1cm, y=1cm, line cap=round, line join=round]
  % 开始
  \draw[rounded corners=6pt] (-0.95,0.42) rectangle (0.95,-0.42);
  \node at (0,0) {开始};
  \draw[->] (0,-0.42) -- (0,-0.87);
  % 输入 x
  \draw (-1.45,-0.87) -- (1.45,-0.87) -- (1.25,-1.72) -- (-1.65,-1.72) -- cycle;
  \node at (0,-1.30) {输入 $x$};
  \draw[->] (0,-1.72) -- (0,-2.16);
  % n=0
  \draw (-1.10,-2.16) rectangle (1.10,-2.96);
  \node at (0,-2.56) {$n=0$};
  \draw[->] (0,-2.96) -- (0,-3.45);
  % 判断框
  \draw (0,-3.45) -- (2.55,-4.55) -- (0,-5.65) -- (-2.55,-4.55) -- cycle;
  \node at (0,-4.55) {$x^2-4x+3\leqslant 0$};
  % 是：向下进入循环体
  \draw[->] (0,-5.65) -- (0,-6.10);
  \node at (0.45,-5.90) {是};
  \draw (-1.10,-6.10) rectangle (1.10,-6.90);
  \node at (0,-6.50) {$x=x+1$};
  \draw[->] (0,-6.90) -- (0,-7.35);
  \draw (-1.10,-7.35) rectangle (1.10,-8.15);
  \node at (0,-7.75) {$n=n+1$};
  % 回边：向下、向左、向上，再水平进入判断框左顶点
  \draw[->] (-1.10,-7.75) -- (-3.45,-7.75) -- (-3.45,-4.55) -- (-2.55,-4.55);
  % 否：向右输出
  \draw[->] (2.55,-4.55) -- (3.95,-4.55);
  \node at (3.20,-4.20) {否};
  \draw (3.95,-4.10) -- (6.35,-4.10) -- (6.15,-5.00) -- (3.75,-5.00) -- cycle;
  \node at (5.05,-4.55) {输出 $n$};
  \draw[->] (5.05,-5.00) -- (5.05,-5.45);
  % 结束
  \draw[rounded corners=6pt] (4.10,-5.45) rectangle (6.00,-6.29);
  \node at (5.05,-5.87) {结束};
\end{tikzpicture}
\end{center}
\nopagebreak[4]
\end{question}
\begin{question}
\gkti{12}在 $\triangle ABC$ 中，已知 $\overrightarrow{AB}\cdot\overrightarrow{AC}=\tan A$，当 $A=\dfrac\pi6$ 时，$\triangle ABC$ 的面积为\gkblank{}。
\end{question}
\begin{question}
\gkti{13}三棱锥 $P-ABC$ 中，$D$，$E$ 分别为 $PB$，$PC$ 的中点，记三棱锥 $D-ABE$ 的体积为 $V_1$，$P-ABC$ 的体积为 $V_2$，则 $\dfrac{V_1}{V_2}=$\gkblank{}。
\end{question}
\begin{question}
\gkti{14}若 $\left(ax^2+\dfrac bx\right)^6$ 的展开式中 $x^3$ 项的系数为 $20$，则 $a^2+b^2$ 的最小值为\gkblank{}。
\end{question}
\begin{question}
\gkti{15}已知函数 $y=f(x)$（$x\in\mathbf R$）。对函数 $y=g(x)$（$x\in I$），定义 $g(x)$ 关于 $f(x)$ 的“对称函数”为函数 $y=h(x)$（$x\in I$），$y=h(x)$ 满足：对任意 $x\in I$，两个点 $(x,h(x))$，$(x,g(x))$ 关于点 $(x,f(x))$ 对称。若 $h(x)$ 是 $g(x)=\sqrt{4-x^2}$ 关于 $f(x)=3x+b$ 的“对称函数”，且 $h(x)>g(x)$ 恒成立，则实数 $b$ 的取值范围是\gkblank{}。
\end{question}
\gkpart{三、解答题}
\begin{question}
\gkti{16}已知向量 $\boldsymbol a=(m,\cos2x)$，$\boldsymbol b=(\sin2x,n)$，函数 $f(x)=\boldsymbol a\cdot\boldsymbol b$，且 $y=f(x)$ 的图象过点 $\left(\dfrac\pi{12},\sqrt3\right)$ 和点 $\left(\dfrac{2\pi}3,-2\right)$。
（1） 求 $m$，$n$ 的值；
（2） 将 $y=f(x)$ 的图象向左平移 $\varphi$（$0<\varphi<\pi$）个单位后得到函数 $y=g(x)$ 的图象。若 $y=g(x)$ 的图象上各最高点到点 $(0,3)$ 的距离的最小值为 $1$，求 $y=g(x)$ 的单调增区间。
\Needspace{0.4\textheight}
\end{question}
\begin{question}
\gkti{17}如图，在四棱柱 $ABCD-A_1B_1C_1D_1$ 中，底面 $ABCD$ 是等腰梯形，$\angle DAB=60^\circ$，$AB=2CD=2$，$M$ 是线段 $AB$ 的中点。
\nopagebreak[4]
\begin{center}% ============================================================
% f2014_shandong_science_q17_prism.tex
% 2014 山东卷（理科）第 17 题：四棱柱 ABCD-A1B1C1D1
% 来源 PDF：普通高考/2014/2014山东理.pdf
% 原页：第 1 页（A3 横排）右栏，印面 907，第 17 题题干下方
% 结构：底面 ABCD 为等腰梯形，AB=2CD=2，∠DAB=60°，M 为 AB 中点。
%       实线：AB、AA1、A1D1、D1C1、C1B1、BB1、A1B1；
%       虚线（不可见棱）：AD、DC、CB、DD1、CC1；
%       虚线（辅助线）：D1A、D1M、C1M。
% 比例：AB:CD=2:1 按屏幕长 4.6:2.3 画出，侧棱方向取 (0.95,2.85)。
% ============================================================
\begin{tikzpicture}[>=Stealth, line width=0.7pt, font=\small,
  x=1cm, y=1cm, line cap=round, line join=round]
  % 底面：AB=4.6（屏幕长），CD=2.3，梯形
  \coordinate (A)  at (0.10,0.30);
  \coordinate (B)  at (4.70,0.30);
  \coordinate (M)  at (2.40,0.30);
  \coordinate (D)  at (1.45,2.05);
  \coordinate (C)  at (3.75,2.05);
  % 侧棱平移向量 (0.95,2.85)
  \coordinate (A1) at (1.05,3.15);
  \coordinate (B1) at (5.65,3.15);
  \coordinate (D1) at (2.40,4.90);
  \coordinate (C1) at (4.70,4.90);
  % 不可见棱
  \draw[dashed] (A) -- (D);
  \draw[dashed] (D) -- (C);
  \draw[dashed] (C) -- (B);
  \draw[dashed] (D) -- (D1);
  \draw[dashed] (C) -- (C1);
  % 辅助线
  \draw[dashed] (D1) -- (A);
  \draw[dashed] (D1) -- (M);
  \draw[dashed] (C1) -- (M);
  % 可见棱
  \draw (A) -- (B);
  \draw (A) -- (A1);
  \draw (A1) -- (D1);
  \draw (D1) -- (C1);
  \draw (C1) -- (B1);
  \draw (B) -- (B1);
  \draw (A1) -- (B1);
  % 顶点与中点标注
  \node[below left]  at (A)  {$A$};
  \node[below right] at (B)  {$B$};
  \node[below]       at (M)  {$M$};
  \node at (1.10,1.94) {$D$};
  \node at (3.98,1.84) {$C$};
  \node[left]        at (A1) {$A_1$};
  \node[right]       at (B1) {$B_1$};
  \node[above left]  at (D1) {$D_1$};
  \node[above right] at (C1) {$C_1$};
\end{tikzpicture}
\end{center}
\nopagebreak[4]
（1） 求证：$C_1M\parallel$ 平面 $A_1ADD_1$；
（2） 若 $CD_1$ 垂直于平面 $ABCD$ 且 $CD_1=\sqrt3$，求平面 $C_1D_1M$ 和平面 $ABCD$ 所成的角（锐角）的余弦值。
\Needspace{0.4\textheight}
\end{question}
\begin{question}
\gkti{18}乒乓球台面被球网分隔成甲、乙两部分，如图，甲上有两个不相交的区域 $A$，$B$，乙被划分为两个不相交的区域 $C$，$D$。某次测试要求队员接到落点在甲上的来球后向乙回球。规定：回球一次，落点在 $C$ 上记 $3$ 分，在 $D$ 上记 $1$ 分，其他情况记 $0$ 分。对落点在 $A$ 上的来球，队员小明回球的落点在 $C$ 上的概率为 $\dfrac12$，在 $D$ 上的概率为 $\dfrac13$；对落点在 $B$ 上的来球，小明回球的落点在 $C$ 上的概率为 $\dfrac15$，在 $D$ 上的概率为 $\dfrac35$。假设共有两次来球且落在 $A$，$B$ 上各一次，小明的两次回球互不影响。求：
\nopagebreak[4]
\begin{center}% ============================================================
% f2014_shandong_science_q18_table.tex
% 2014 山东卷（理科）第 18 题：乒乓球台面分区示意图
% 来源 PDF：普通高考/2014/2014山东理.pdf
% 原页：第 2 页（A3 横排）左栏，印面 908，第 18 题题干下方
% 结构：球台台面以斜视角画成平行四边形，中间一条球网（带网带的窄条与
%       露出台面的网顶）把台面分成两部分：左侧为乙区（含不相交的区域
%       C、D，原图以浅灰底纹标示），右侧为甲区（含不相交的区域 A、B）。
% ============================================================
\begin{tikzpicture}[>=Stealth, line width=0.7pt, font=\small,
  x=1cm, y=1cm, line cap=round, line join=round]
  % 台面四角：近边在下，远边在上，两端向右后方倾斜
  \coordinate (NL) at (0.60,0.80);
  \coordinate (NR) at (7.40,0.80);
  \coordinate (FR) at (8.60,2.70);
  \coordinate (FL) at (1.80,2.70);
  % 球网两条网脚线与网顶
  \coordinate (B1) at (3.90,0.80);   % 近边左网脚
  \coordinate (B2) at (4.80,0.80);   % 近边右网脚
  \coordinate (T1) at (5.10,2.70);   % 远边左网脚
  \coordinate (T2) at (6.00,2.70);   % 远边右网脚
  \coordinate (U1) at (5.10,3.30);   % 网顶左
  \coordinate (U2) at (6.00,3.30);   % 网顶右
  % 乙区（网左侧）浅灰底纹
  \fill[black!12] (NL) -- (B1) -- (T1) -- (FL) -- cycle;
  % 网带
  \fill[black!30] (B1) -- (T1) -- (T2) -- (B2) -- cycle;
  % 露出台面的网顶
  \fill[black!30] (T1) -- (U1) -- (U2) -- (T2) -- cycle;
  % 台面轮廓
  \draw (NL) -- (NR) -- (FR) -- (FL) -- cycle;
  % 球网
  \draw (B1) -- (T1) -- (U1) -- (U2) -- (T2) -- (B2);
  \draw (T1) -- (T2);
  % 区域标注
  \node at (1.45,1.20) {$C$};
  \node at (3.05,1.95) {$D$};
  \node at (7.05,2.25) {$A$};
  \node at (5.55,1.30) {$B$};
\end{tikzpicture}
\end{center}
\nopagebreak[4]
（1） 小明两次回球的落点中恰有一次的落点在乙上的概率；
（2） 两次回球结束后，小明得分之和 $\xi$ 的分布列与数学期望。
\end{question}
\begin{question}
\gkti{19}已知等差数列 $\{a_n\}$ 的公差为 $2$，前 $n$ 项和为 $S_n$，且 $S_1$，$S_2$，$S_4$ 成等比数列。
（1） 求数列 $\{a_n\}$ 的通项公式；
（2） 令 $b_n=(-1)^{n-1}\dfrac{4n}{a_na_{n+1}}$，求数列 $\{b_n\}$ 的前 $n$ 项和 $T_n$。
\end{question}
\begin{question}
\gkti{20}设函数 $f(x)=\dfrac{\mathrm e^x}{x^2}-k\left(\dfrac2x+\ln x\right)$（$k$ 为常数，$\mathrm e=2.71828\cdots$ 是自然对数的底数）。
（1） 当 $k\leqslant0$ 时，求函数 $f(x)$ 的单调区间；
（2） 若函数 $f(x)$ 在 $(0,2)$ 内存在两个极值点，求 $k$ 的取值范围。
\end{question}
\begin{question}
\gkti{21}已知抛物线 $C:y^2=2px$（$p>0$）的焦点为 $F$，$A$ 为 $C$ 上异于原点的任意一点，过点 $A$ 的直线 $l$ 交 $C$ 于另一点 $B$，交 $x$ 轴的正半轴于点 $D$，且有 $|FA|=|FD|$。当点 $A$ 的横坐标为 $3$ 时，$\triangle ADF$ 为正三角形。
（1） 求 $C$ 的方程；
（2） 若直线 $l_1\parallel l$，且 $l_1$ 和 $C$ 有且只有一个公共点 $E$，
　　① 证明直线 $AE$ 过定点，并求出定点坐标；
　　② $\triangle ABE$ 的面积是否存在最小值？若存在，请求出最小值；若不存在，请说明理由。

\end{question}